% Pacotes
\usepackage[utf8]{inputenc}
\usepackage[T1]{fontenc}
\usepackage[brazil]{babel}
\usepackage{graphicx}
\usepackage{amsmath}
\usepackage{amssymb}
\usepackage{titlesec}
\usepackage{titletoc} % Para formatar listas
\usepackage{helvet}   % Arial
\usepackage{times}    % Times New Roman
\usepackage{tocloft}  % Para formatar sumário, lista de figuras/tabelas
\usepackage{hyperref}
\usepackage{xcolor}

% Comando personalizado para citação no formato "Autor (ano)".
% Usa \textcite quando disponível (biblatex); caso contrário, usa \citeauthor (ano).
% Evita produzir parênteses duplicados porque \textcite já formata corretamente
% no estilo bibliográfico; o fallback monta manualmente "Autor (ano)".
\providecommand{\citeautor}[1]{}
% Use \textcite provided by biblatex to get the format "Autor (ano)".
% This relies on biblatex being loaded (it is in the preamble). If for any
% reason \textcite is unavailable, the user can change this to the fallback
% form, but in our setup \textcite prints the desired output reliably.
\renewcommand{\citeautor}[1]{\textcite{#1}}
% ===== Fonte padrão do texto =====
\renewcommand{\rmdefault}{ptm} % Times New Roman

% ===== Formatação dos Capítulos =====
\titleformat{\chapter}[hang]
  {\normalfont\bfseries\fontsize{12}{14}\selectfont}
  {\thechapter.}{1em}{}

\titlespacing*{\chapter}{0pt}{-10pt}{20pt}

% ===== Formatação dos títulos das listas (Sumário, Figuras, Tabelas) =====
\renewcommand{\cfttoctitlefont}{\fontfamily{phv}\selectfont\fontsize{16}{18}\bfseries} % Sumário
\renewcommand{\cftloftitlefont}{\fontfamily{phv}\selectfont\fontsize{16}{18}\bfseries} % Lista de figuras
\renewcommand{\cftlottitlefont}{\fontfamily{phv}\selectfont\fontsize{16}{18}\bfseries} % Lista de tabelas


% ===== Espaçamento =====
\linespread{1.5}